\subsection{导子的复合}

我们在本号中假设第 2 号的情形（I）成立，即 A、A'、A'' 是三个类型为 Δ 的分次 K-模，并且给定一个 K-双线性映射 μ: A × A' → A''，它对应于一个次数为 0 的分次 K-线性映射 A ⊗K A' → A''。A ⊕ A' ⊕ A'' 的分次自同态 f，若满足 f(A) ⊂ A、f(A') ⊂ A' 且 f(A'') ⊂ A''，则构成分次结合代数 EndgrK(A ⊕ A' ⊕ A'')（§ 3，第 1 号，例 2）的一个分次子代数。特别地，两个ε-导子（关于(A, A', A'')）可以复合，但不应认为两个ε-导子的复合仍是ε-导子。

在每一个类型为 \(\Delta\) 的分次代数 B 上， \(\varepsilon\) -括号（当 \(\varepsilon = 1\) 时简称括号）由公式 (10) 对两个齐次元素 \(u, v\) 定义：

\[
[ u, v ] _ {\varepsilon} = u v - \varepsilon_ {\deg u, \deg v} v u (\text { denoted   simply   by } [ u, v ] \text { if } \varepsilon = 1).
\]

将此映射线性扩张，就得到一个从 B × B 到 B 的 K-双线性映射 \((u, v) \mapsto [u, v]_{\varepsilon}\) 。于是，对 B 中齐次的 u 和 v

\[
[ v, u ] _ {\varepsilon} = - \varepsilon_ {\deg u, \deg v} [ u, v ] _ {\varepsilon}.
\]

将此定义应用于分次代数 \(\operatorname{Endgr}_{\mathbf{K}}(\mathbf{A} \oplus \mathbf{A}' \oplus \mathbf{A}'')\) ，两个分次自同态的 \(\varepsilon\) -括号便由此得到定义。

命题 1. 设 \(d_1, d_2\) 是两个 \(\varepsilon\) -导子（属于 \((\mathbf{A}, \mathbf{A}', \mathbf{A} ' )\) ，次数分别为 \(\delta_1, \delta_2\) ）。则 \(\varepsilon\) -括号

\[
[ d _ {1}, d _ {2} ] _ {\varepsilon} = d _ {1} \circ d _ {2} - \varepsilon_ {\delta_ {1}, \delta_ {2}} d _ {2} \circ d _ {1}
\]

是一个 \(\varepsilon\) -导子（次数为 \(\delta_1 + \delta_2\) ）。此外，若 \(d\) 是一个 \(\varepsilon\) -导子（属于 \((\mathbf{A}, \mathbf{A}', \mathbf{A}'')\) ，次数为 \(\delta\) ），且 \(\varepsilon_{\delta, \delta} = -1\) ，则 \(d^2 = d \circ d\) 是一个导子。

设 \(x \in A\) 是次数为 \(\xi\) 的齐次元素；对所有 \(y \in A'\) ，

\[
\begin{array}{r l} d _ {1} (d _ {2} (x y)) & = ((d _ {1} d _ {2}) (x)) y + \varepsilon_ {\delta_ {1}, \delta_ {2} + \xi} (d _ {2} x) (d _ {1} y) \\ & \quad + \varepsilon_ {\delta_ {2}, \xi} (d _ {1} x) (d _ {2} y) + \varepsilon_ {\delta_ {1} + \delta_ {2}, \xi} x ((d _ {1} d _ {2}) (y)) \end{array}\tag{11}
\]

这里利用了第 1 号的公式 (1)。交换 \(d_1\) 和 \(d_2\) 的角色，并再次利用 (1)（第 1 号）化简后，我们得到

\[
\begin{array}{r l} (d _ {1} d _ {2}) (x y) - \varepsilon_ {\delta_ {1}, \delta_ {2}} (d _ {2} d _ {1}) (x y) = & ((d _ {1} d _ {2}) (x)) y - \varepsilon_ {\delta_ {1}, \delta_ {2}} ((d _ {2} d _ {1}) (x)) y \\ & + \varepsilon_ {\delta_ {1} + \delta_ {2}, \xi} x ((d _ {1} d _ {2}) (y)) \\ & - \varepsilon_ {\delta_ {1}, \delta_ {2}} \varepsilon_ {\delta_ {1} + \delta_ {2}, \xi} x ((d _ {2} d _ {1}) (y)) \end{array}
\]

即，记 \(d = [d_1, d_2]_{\varepsilon}\) 和 \(\delta = \delta_1 + \delta_2\) ,

\[
d (x y) = (d x) y + \varepsilon_ {\delta , \xi} x (d y)
\]

这证明了 \(d\) 是一个 \(\varepsilon\) -导子。

另一方面，若在 (11) 中令 \(d_1 = d_2 = d\) ， \(\delta_1 = \delta_2 = \delta\) 以及 \(\varepsilon_{\delta,\delta} = -1\) ，则由 (1)，此时 \(\varepsilon_{\delta,\delta + \xi} = -\varepsilon_{\delta,\xi}\) ，故得到

\[
d ^ {2} (x y) = \left(d ^ {2} x\right) y + \varepsilon_ {2 5, \xi} x \left(d ^ {2} y\right)
\]

又由于 \(\varepsilon_{2\delta, \xi} = 1\) ，可以看出 \(d^2\) 是一个导子。

推论。假设 \(\Delta = \mathbf{Z}\) 。那么：

\begin{enumerate}
\def\labelenumi{(\roman{enumi})}
\item
  一个反导子的平方是一个导子。
\item
  两个导子之间的括号运算仍是一个导子。
\item
  反导子与偶数次导子的括号运算仍为反导子。
\item
  若 \(d_1\) 和 \(d_2\) 是反导子，则 \(d_1d_2 + d_2d_1\) 是一个导子。
\end{enumerate}

在本号开头的假设下，现在考虑一个有限序列 \(\mathbf{D} = (d_i)_{1 \leqslant i \leqslant n}\) ，它由 \((\mathbf{A}, \mathbf{A}', \mathbf{A}'')\) 的两两可交换的导子组成。对每个多项式 \(\mathrm{P}(\mathbf{X}_1, \ldots, \mathbf{X}_n)\) （属于代数 \(\mathbf{K}[\mathbf{X}_1, \ldots, \mathbf{X}_n]\) ），元素 \(\mathrm{P}(d_1, \ldots, d_n)\) （ \(\operatorname{Endgr}_{\mathbf{K}}(\mathbf{A} \oplus \mathbf{A}' \oplus \mathbf{A}'')\) 中的）便有定义（§ 2，第 9 号）；其简写记号为 \(\mathrm{P}(\mathbf{D})\) 。

命题 2. 在上述假设与记号下，考虑 \(2n\) 个未定元 \(\mathrm{T}_1, \ldots, \mathrm{T}_n\) ， \(\mathrm{T}_1', \ldots, \mathrm{T}_n'\) ，并且对每个多项式 \(\mathbf{F} \in \mathbf{K}[\mathbf{X}_1, \ldots, \mathbf{X}_n]\) ，记 \(\mathbf{F}(\mathbf{T}) = \mathbf{F}(\mathbf{T}_1, \ldots, \mathbf{T}_n)\) ， \(\mathbf{F}(\mathbf{T}') = \mathbf{F}(\mathbf{T}_1', \ldots, \mathbf{T}_n')\) 以及

\[
\mathrm{F} (\mathbf {T} + \mathbf {T} ^ {\prime}) = \mathrm{F} (\mathrm{T} _ {1} + \mathrm{T} _ {1} ^ {\prime}, \dots , \mathrm{T} _ {n} + \mathrm{T} _ {n} ^ {\prime}).
\]

假设

\[
\mathbf {P} (\mathbf {T} + \mathbf {T} ^ {\prime}) = \sum_ {i} \mathbf {Q} _ {i} (\mathbf {T}) \mathbf {R} _ {i} (\mathbf {T} ^ {\prime})
\]

其中 \(\mathbf{Q}_i\) 和 \(\mathbf{R}_i\) 属于 \(\mathbf{K}[\mathbf{X}_1,\ldots ,\mathbf{X}_n]\) 。那么，对于所有 \(x\in \mathbf{A}\) 和 \(y\in \mathbf{A}'\) ,

\[
\mathbf {P} (\mathbf {D}) (x y) = \sum_ {i} \left(\mathbf {Q} _ {i} (\mathbf {D}) x\right) (\mathbf {R} _ {i} (\mathbf {D}) y).\tag{12}
\]

我们引入另外 \(n\) 个未定元 \(T_1'', \ldots, T_n''\) ，并考虑多项式代数 \(K[T_1, \ldots, T_n, T_1', \ldots, T_n', T_1'', \ldots, T_n''] = B\) ；另一方面，我们考虑 K-模 \(M\) ，它由 \(A \times A'\) 到 \(A''\) 的双线性映射组成；在 \(M\) 上定义一个 B-模结构如下：对每个 K-双线性映射 \(f \in M\) 以及 \(1 \leqslant i \leqslant n\) ，

\[
\left\{ \begin{array}{l} (\mathrm{T} _ {i} f) (a, a ^ {\prime}) = f (d _ {i} a, a ^ {\prime}) \\ (\mathrm{T} _ {i} ^ {\prime} f) (a, a ^ {\prime}) = f (a, d _ {i} a ^ {\prime}) \\ (\mathrm{T} _ {i} ^ {\prime \prime} f) (a, a ^ {\prime}) = d _ {i} (f (a, a ^ {\prime})). \end{array} \right.\tag{13}
\]

由于诸 \(d_{i}\) 两两可交换，可见对每个多项式 \(\mathbf{F} \in \mathbf{K}[\mathbf{X}_{1}, \ldots, \mathbf{X}_{n}]\) ，有 \((\mathbf{F}(\mathbf{T})f)(a, a') = f(\mathbf{F}(\mathbf{D})a, a')\) ，

\[
(\mathbf {F} (\mathbf {T} ^ {\prime}) f) (a, a ^ {\prime}) = f (a, \mathbf {F} (\mathbf {D}) a ^ {\prime})
\]

以及 \((\mathbf{F}(\mathbf{T}^{\prime \prime})f) = \mathbf{F}(\mathbf{D})(f(a,a^{\prime}))\) 。因此，为证明 (12)，只需证明

\[
(\mathbf {P} (\mathbf {T} ^ {\prime \prime}) - \sum_ {i} \mathbf {Q} _ {i} (\mathbf {T}) \mathbf {R} _ {i} (\mathbf {T} ^ {\prime})) \cdot \mu = 0\tag{14}
\]

或等价地 \((\mathrm{P}(\mathbf{T}^{\prime\prime}) - \mathrm{P}(\mathbf{T} + \mathbf{T}^{\prime}))\) . \(\mu = 0\) 在 B-模 M 中成立。现在，诸 \(d_{i}\) 是导子这一假设也可以表述为：对 \(1 \leqslant i \leqslant n\) ，

\[
\left(\mathrm{T} _ {i} ^ {\prime \prime} - \mathrm{T} _ {i} - \mathrm{T} _ {i} ^ {\prime}\right). \mu = 0\tag{15}
\]

在 B-模 M 中成立。依次考虑多项式

\[
\mathrm{P} \left(\mathrm{T} _ {1} ^ {\prime \prime}, \mathrm{T} _ {2} ^ {\prime \prime}, \dots , \mathrm{T} _ {n} ^ {\prime \prime}\right) - \mathrm{P} \left(\mathrm{T} _ {1} + \mathrm{T} _ {1} ^ {\prime}, \mathrm{T} _ {2} ^ {\prime \prime}, \dots , \mathrm{T} _ {n} ^ {\prime \prime}\right)
\]

\[
\mathrm{P} \left(\mathrm{T} _ {1} + \mathrm{T} _ {1} ^ {\prime}, \mathrm{T} _ {2} ^ {\prime \prime}, \dots , \mathrm{T} _ {n} ^ {\prime \prime}\right) - \mathrm{P} \left(\mathrm{T} _ {1} + \mathrm{T} _ {1} ^ {\prime}, \mathrm{T} _ {2} + \mathrm{T} _ {2} ^ {\prime}, \dots , \mathrm{T} _ {n} ^ {\prime \prime}\right)
\]

\[
\mathrm{P} \left(\mathrm{T} _ {1} + \mathrm{T} _ {1} ^ {\prime}, \dots , \mathrm{T} _ {n - 1} + \mathrm{T} _ {n - 1} ^ {\prime}, \mathrm{T} _ {n} ^ {\prime \prime}\right) - \mathrm{P} \left(\mathrm{T} _ {1} + \mathrm{T} _ {1} ^ {\prime}, \dots , \mathrm{T} _ {n - 1} + \mathrm{T} _ {n - 1} ^ {\prime}, \mathrm{T} _ {n} + \mathrm{T} _ {n} ^ {\prime}\right)
\]

就可以看出，差 \(\mathbf{P}(\mathbf{T}^{\prime \prime}) - \mathbf{P}(\mathbf{T} + \mathbf{T}^{\prime})\) 可以写成如下形式

\[
\sum_ {i = 1} ^ {n} \left(\mathbf {T} _ {i} ^ {\prime \prime} - \mathbf {T} _ {i} - \mathbf {T} _ {i} ^ {\prime}\right) \mathrm{G} _ {i} (\mathbf {T}, \mathbf {T} ^ {\prime}, \mathbf {T} ^ {\prime \prime})
\]

其中 \(G_{i}\) 是 B 的元素。因此关系 (14) 是关系 (15) 的直接推论。

推论（莱布尼茨公式）。设 \(d_{i}\) ( \(1 \leqslant i \leqslant n\) ) 是 \(n\) 个导子（属于 \((\mathbf{A}, \mathbf{A}', \mathbf{A}'' )\) ，两两可交换）。对所有 \(\alpha = (\alpha_{1}, \ldots, \alpha_{n}) \in \mathbf{N}^{n}\) ，我们记

\[
d ^ {\alpha} = d _ {1} ^ {\alpha_ {1}} d _ {2} ^ {\alpha_ {2}} \dots d _ {n} ^ {\alpha_ {n}}.\tag{16}
\]

那么，对 \(x \in \mathbf{A}\) 和 \(y \in \mathbf{A}'\) ，

\[
d ^ {\alpha} (x y) = \sum_ {\beta + \gamma = \alpha} ((\beta , \gamma)) d ^ {\beta} (x) d ^ {\gamma} (y)\tag{17}
\]

其中我们写下了（采用本章开头引入的记号）

\[
((\beta , \gamma)) = (\beta + \gamma)! / (\beta ! \gamma!).\tag{18}
\]

这由多项式公式（I，§ 8，第 2 号）

\[
(\mathbf {T} + \mathbf {T} ^ {\prime}) ^ {\alpha} = \sum_ {\beta + \gamma = \alpha} ((\beta , \gamma)) \mathbf {T} ^ {\beta} \mathbf {T} ^ {\prime \gamma}
\]

以及命题 2 立即得出。

